\chapter{VLA 연구를 읽는 13개 질문}

\section{장의 목표}

앞의 두 장에서는 VLA를 언어 조건부 폐루프 정책으로 정의하고, 관측 인코더, 언어 인코더, 멀티모달 결합, 행동 표현, 정책 head, 실행 loop로 구조를 분해했다. 이제 필요한 것은 논문을 읽을 때 같은 기준으로 비교하는 방법이다. VLA 논문은 빠르게 늘고 있으며, 제목만 보면 모두 ``generalist robot policy'', ``efficient VLA'', ``reasoning VLA''처럼 보일 수 있다. 그러나 실제 기여는 서로 다르다. 어떤 논문은 행동 표현을 바꾸고, 어떤 논문은 fine-tuning recipe를 바꾸며, 어떤 논문은 벤치마크를 만든다.

이 장의 목적은 VLA 논문을 읽는 13개 질문을 공학적 평가 도구로 정식화하는 것이다. 여기서 13개 질문은 요약 양식이 아니다. 논문이 실제로 무엇을 주장하며, 그 주장이 어떤 실험으로 검증되었고, 어디까지 믿을 수 있는지를 확인하는 audit checklist이다.

\section{13개 질문의 전체 구조}

이 책에서 사용하는 13개 질문은 다음과 같다.

\begin{enumerate}
  \item 배경: 이 연구는 어떤 분야와 임무 맥락에 놓이는가?
  \item 문제: 구체적으로 어떤 기능 또는 과제를 풀려 하는가?
  \item 기존 한계: 이전 접근의 어떤 병목 때문에 새 방법이 필요한가?
  \item 목표: 논문이 달성하려는 목표는 측정 가능한가?
  \item 방법: 제안 시스템은 어떤 입력, 출력, 모듈, 학습 절차를 갖는가?
  \item 핵심 아이디어: 기술적으로 무엇이 새롭거나 중요한가?
  \item 검증: 어떤 데이터셋, 시뮬레이터, 실제 로봇, 시나리오로 평가했는가?
  \item 결과: 어떤 지표에서 얼마나 좋아졌는가?
  \item 비교: baseline은 공정하고 충분히 강한가?
  \item 의의: 결과가 VLA와 로봇 학습에서 무엇을 가능하게 하는가?
  \item 한계: 남은 제약과 실패 모드는 무엇인가?
  \item 향후 과제: 다음 연구 단계는 구체적으로 무엇인가?
  \item 자원 공개: 코드, 데이터, checkpoint, benchmark가 공개되었는가?
\end{enumerate}

이 순서는 중요하다. 많은 논문 요약은 방법과 결과부터 시작하지만, 그렇게 읽으면 왜 그 방법이 필요한지 놓치기 쉽다. 반대로 배경과 문제만 길게 읽으면 논문이 실제로 무엇을 검증했는지 흐려진다. 13개 질문은 논문의 주장 흐름을 배경에서 공개 자원까지 닫힌 형태로 추적한다.

\section{질문 1--4: 문제 설정을 검증하기}

첫 네 질문은 논문의 문제 설정을 검사한다. 좋은 VLA 논문은 배경, 문제, 기존 한계, 목표가 서로 직접 연결된다. 예를 들어 ``VLA가 느리다''는 배경이 있다면 문제는 추론 latency여야 하고, 기존 한계는 autoregressive decoding 또는 큰 backbone의 계산량이어야 하며, 목표는 latency를 줄이면서 success rate를 유지하는 것이어야 한다. 만약 배경은 real-world deployment인데 실험은 작은 offline dataset에만 머문다면 문제 설정과 검증이 분리된 것이다.

\subsection{배경}

배경은 연구의 큰 위치를 정한다. VLA 논문에서 배경은 보통 다음 중 하나이다.

\begin{itemize}
  \item manipulation: 언어 지시를 실제 조작 행동으로 바꾸는 문제
  \item efficient VLA: 큰 VLA를 실제 제어 주기에 맞추는 문제
  \item action representation: 연속 행동을 모델이 다루기 쉬운 표현으로 바꾸는 문제
  \item benchmark: 모델 성능을 공정하게 비교하는 문제
  \item reasoning and planning: 장기 과제와 단계적 의사결정을 다루는 문제
  \item robustness and safety: 분포 변화와 실패 모드를 다루는 문제
\end{itemize}

배경을 쓸 때는 ``로봇이 중요하다'' 같은 일반론으로 끝내면 안 된다. 어떤 로봇 과제, 어떤 데이터 흐름, 어떤 deployment constraint가 문제인지 지정해야 한다.

\subsection{문제}

문제는 구체적이어야 한다. ``VLA의 성능을 높인다''는 문제 정의가 아니다. 다음처럼 변수와 평가 기준으로 바꿀 수 있어야 한다.

\begin{equation}
  \max_\theta \; \mathrm{SR}(\theta)
  \quad
  \mathrm{s.t.}
  \quad
  \tau_{\mathrm{model}}(\theta) \le \tau_{\max}.
  \label{eq:problem_efficiency}
\end{equation}

이 식은 efficient VLA 논문의 문제 정의에 가깝다. 성공률을 최대화하되 모델 추론 시간이 허용 범위를 넘지 않아야 한다. Action representation 논문이라면 문제는 quantization error와 policy loss 사이의 균형이 될 수 있다.

\begin{equation}
  \min_{Q,\theta}
  \mathcal{L}_{\mathrm{policy}}(\theta; Q)
  +
  \lambda \epsilon_Q.
  \label{eq:problem_tokenization}
\end{equation}

문제 정의가 이런 식으로 쓰이지 않더라도, 독자는 논문을 읽으며 암묵적 objective를 복원해야 한다.

\subsection{기존 한계}

기존 한계는 논문의 필요성을 만드는 부분이다. 약한 논문은 기존 연구를 단순히 ``비싸다'', ``느리다'', ``일반화가 약하다''라고 말하지만, 어떤 구조적 원인 때문에 그런지 설명하지 않는다. 강한 논문은 한계를 모듈 수준으로 지정한다. 예를 들어 action tokenization이 약하다면 codebook 설계가 문제인지, temporal consistency가 문제인지, high-frequency control과 token prediction의 mismatch가 문제인지 밝혀야 한다.

\subsection{목표}

목표는 측정 가능해야 한다. 목표가 ``더 일반적인 VLA''라면 어떤 test split에서 일반화를 측정하는지 필요하다. 목표가 ``real-time VLA''라면 loop frequency와 latency budget이 필요하다. 목표가 ``robust VLA''라면 어떤 perturbation, distractor, unseen object, unseen instruction을 넣는지 필요하다. 목표가 측정 불가능하면 실험이 좋아 보여도 주장이 닫히지 않는다.

\section{질문 5--6: 방법의 기술적 중심 찾기}

질문 5와 6은 방법을 읽는 단계이다. 여기서 중요한 것은 방법 설명과 핵심 아이디어를 분리하는 것이다. 방법은 시스템 전체의 구성이고, 핵심 아이디어는 그중 논문이 실제로 새로 만든 부분이다.

\subsection{방법}

VLA 방법 설명은 최소한 다음 다섯 항목을 포함해야 한다.

\begin{enumerate}
  \item 입력: 이미지, 비디오, proprioception, force, 언어 지시
  \item 출력: action token, end-effector command, joint command, trajectory, action chunk
  \item backbone: VLM, LLM, transformer, diffusion/flow model
  \item 학습: behavior cloning, instruction tuning, reinforcement fine-tuning, adapter tuning
  \item 추론: closed-loop frequency, chunk length, decoding step, safety filter
\end{enumerate}

논문이 모델 그림만 보여 주고 위 항목을 명확히 하지 않으면 재현성이 약하다. 특히 action space와 controller interface는 반드시 확인해야 한다.

\subsection{핵심 아이디어}

핵심 아이디어는 한 문장으로 줄일 수 있어야 한다. 예를 들어 다음과 같이 쓸 수 있다.

\begin{itemize}
  \item Action tokenization 논문: 연속 로봇 행동을 VLM/LLM이 예측하기 쉬운 discrete token sequence로 바꾼다.
  \item Efficient VLA 논문: 모든 token을 같은 계산량으로 처리하지 않고 task-relevant token만 선택적으로 계산한다.
  \item Post-training 논문: pretrained VLA를 simulator reward 또는 verified feedback으로 다시 보정한다.
  \item Benchmark 논문: 다양한 과제와 domain randomization으로 VLA의 robustness를 같은 조건에서 비교한다.
\end{itemize}

핵심 아이디어가 ``새로운 framework를 제안한다''로만 남으면 부족하다. 어떤 병목을 어떤 mechanism으로 줄이는지가 보여야 한다.

\section{질문 7--9: 검증의 공정성 보기}

질문 7--9는 논문의 실험을 평가한다. VLA 분야에서 실험은 특히 조심해서 봐야 한다. 로봇 플랫폼, reset 방식, instruction template, object set, camera view, demonstration quality가 조금만 달라도 성공률이 달라진다.

\subsection{검증}

검증은 다음 세 수준으로 나눌 수 있다.

\begin{table}[t]
  \centering
  \caption{VLA 검증 수준}
  \label{tab:validation_levels}
  \begin{tabularx}{\textwidth}{>{\bfseries}p{32mm}X}
    \toprule
    수준 & 의미 \\
    \midrule
    Offline dataset & 저장된 trajectory와 observation에서 action prediction 성능을 평가한다. \\
    Simulation benchmark & 시뮬레이터에서 policy를 실행해 success rate와 robustness를 평가한다. \\
    Real robot & 실제 센서, controller, latency, 물리 접촉을 포함해 task success를 평가한다. \\
    \bottomrule
  \end{tabularx}
\end{table}

세 수준은 서로 대체되지 않는다. offline action prediction이 좋아도 closed-loop success가 낮을 수 있고, simulation success가 좋아도 real robot에서는 접촉과 센서 노이즈 때문에 실패할 수 있다. 논문이 어떤 수준의 검증을 했는지 먼저 분리해야 한다.

\subsection{결과}

결과는 숫자로 읽되, 숫자의 조건을 함께 읽어야 한다. 성공률은 다음처럼 조건부 지표이다.

\begin{equation}
  \mathrm{SR}
  =
  \mathrm{SR}
  (
  \mathcal{T},
  \mathcal{E},
  \mathcal{O},
  \mathcal{L},
  \mathcal{R}
  ),
\end{equation}

여기서 $\mathcal{T}$는 task set, $\mathcal{E}$는 environment distribution, $\mathcal{O}$는 object set, $\mathcal{L}$은 language distribution, $\mathcal{R}$은 robot platform이다. 단일 숫자로 쓰인 success rate는 이 조건들을 압축한 결과일 뿐이다.

Latency도 마찬가지이다. 모델 forward time만 보고하면 실제 deployment latency를 알 수 없다. end-to-end latency는 센서 입력, 전처리, 모델 추론, decoding, safety filter, controller interface를 포함해야 한다.

\subsection{비교}

비교에서 가장 중요한 것은 baseline의 강도이다. 약한 baseline만 이기면 논문 주장이 과장될 수 있다. VLA 논문에서 baseline은 다음 그룹을 고려해야 한다.

\begin{itemize}
  \item 같은 데이터로 학습한 기존 VLA
  \item 같은 action space를 쓰는 diffusion policy 또는 visuomotor policy
  \item frozen VLM + planner + controller pipeline
  \item ablation: backbone 고정, action head 변경, data size 변경, chunk length 변경
\end{itemize}

특히 action representation 논문은 같은 backbone과 같은 데이터에서 action representation만 바꾸는 ablation이 필요하다. Efficient VLA 논문은 같은 success rate에서 latency가 줄었는지, 같은 latency에서 success rate가 올랐는지 분리해서 보여야 한다.

\section{질문 10--12: 의미와 한계를 분리하기}

질문 10--12는 논문 결과를 해석하는 단계이다. 여기서 흔한 오류는 의의를 과장하거나 한계를 future work 한 줄로 처리하는 것이다.

\subsection{의의}

의의는 결과가 가능하게 하는 것을 말한다. VLA 논문에서 의의는 보통 다음 중 하나이다.

\begin{enumerate}
  \item 더 작은 GPU 또는 더 빠른 제어 주기에서 VLA를 실행할 수 있다.
  \item 더 다양한 조작 과제에서 language-conditioned policy를 쓸 수 있다.
  \item action representation을 바꾸어 VLM과 로봇 제어 사이의 mismatch를 줄인다.
  \item benchmark를 통해 모델 간 비교와 failure diagnosis가 가능해진다.
  \item reasoning, memory, world model을 통해 long-horizon task로 확장한다.
\end{enumerate}

의의는 논문이 실제로 검증한 범위를 넘으면 안 된다. simulation만 했다면 real-world deployment 의의는 제한적으로 써야 한다. offline dataset만 했다면 closed-loop robot policy로서의 의의는 보수적으로 써야 한다.

\subsection{한계}

한계는 두 종류로 나누어야 한다. 첫째는 논문이 명시한 한계이다. 둘째는 독자가 근거로부터 도출하는 한계이다. 예를 들어 논문이 ``실제 로봇 실험은 하지 않았다''고 말하면 명시적 한계이다. 논문이 말하지 않았지만 benchmark가 template instruction만 사용했다면, 자연어 일반화가 제한적이라는 도출 한계를 쓸 수 있다.

\begin{warningbox}{한계 작성 원칙}
한계는 논문을 깎아내리는 문장이 아니라 주장 범위를 정확히 닫는 문장이다. 좋은 한계 분석은 다음 실험을 설계할 수 있게 만든다.
\end{warningbox}

\subsection{향후 과제}

향후 과제는 ``더 연구해야 한다''가 아니어야 한다. VLA 논문에서 의미 있는 향후 과제는 구체적 실험이나 시스템 확장으로 써야 한다.

\begin{itemize}
  \item unseen object와 unseen instruction을 분리한 generalization test
  \item real robot latency와 safety filter를 포함한 deployment benchmark
  \item 실패 trajectory와 recovery data를 포함한 post-training
  \item multi-embodiment action representation의 통일
  \item force/tactile signal을 포함한 contact-rich VLA
\end{itemize}

\section{질문 13: 공개 자원 확인}

마지막 질문은 공개 자원이다. 공개 자원은 단순 부록이 아니다. VLA는 로봇 플랫폼, controller, dataset split, evaluation script가 복잡하기 때문에 공개 여부가 연구의 검증 가능성을 크게 좌우한다.

공개 자원은 다음 수준으로 나눈다.

\begin{table}[t]
  \centering
  \caption{VLA 공개 자원의 수준}
  \label{tab:resource_levels}
  \begin{tabularx}{\textwidth}{>{\bfseries}p{28mm}X}
    \toprule
    수준 & 확인할 내용 \\
    \midrule
    Paper only & 논문과 PDF만 있고 재현 코드는 없다. \\
    Code & 학습 또는 추론 코드가 공개되어 있다. \\
    Checkpoint & pretrained 또는 fine-tuned model weight가 공개되어 있다. \\
    Dataset & trajectory, language instruction, split metadata가 공개되어 있다. \\
    Benchmark & evaluation script, success checker, environment setup이 공개되어 있다. \\
    Full stack & code, checkpoint, data, benchmark, robot/simulator setup이 함께 제공된다. \\
    \bottomrule
  \end{tabularx}
\end{table}

논문 요약에서 ``코드 공개''라고 쓸 때는 실제 URL을 확인해야 한다. GitHub repository가 있어도 empty repository이거나 inference만 있고 training code가 없을 수 있다. Dataset page가 있어도 license가 제한적일 수 있다. 이 책에서는 공개 자원을 평가할 때 링크 존재와 사용 가능성을 구분한다.

\section{13개 질문을 표로 적용하기}

실제 논문을 읽을 때는 다음 표를 채우면 된다.

\begin{table}[t]
  \centering
  \caption{13-question reading template}
  \label{tab:thirteen_template}
  \begin{tabularx}{\textwidth}{>{\bfseries}p{25mm}X}
    \toprule
    항목 & 작성 내용 \\
    \midrule
    배경 & 분야, task family, deployment context \\
    문제 & 해결하려는 기능과 평가 가능한 objective \\
    기존 한계 & 이전 접근의 모듈 수준 병목 \\
    목표 & 측정 가능한 target metric 또는 system property \\
    방법 & 입력, 출력, backbone, loss, inference loop \\
    핵심 아이디어 & contribution mechanism 한 문장 \\
    검증 & dataset, simulator, robot, scenario \\
    결과 & metric과 수치, 조건 \\
    비교 & baseline, ablation, fairness \\
    의의 & 검증 범위 안에서 가능한 것 \\
    한계 & 명시 한계와 도출 한계 \\
    향후 과제 & 다음 실험 또는 시스템 확장 \\
    자원 공개 & code, data, checkpoint, benchmark 링크 \\
    \bottomrule
  \end{tabularx}
\end{table}

이 표는 간단하지만 강력하다. 한 논문을 이 표로 채우지 못하면 보통 두 가지 중 하나이다. 논문 자체가 필요한 정보를 충분히 제공하지 않았거나, 독자가 아직 논문의 핵심을 잡지 못한 것이다.

\section{나쁜 요약과 좋은 요약}

13-question framework는 문장을 길게 만드는 도구가 아니라 근거 없는 요약을 줄이는 도구이다. 예를 들어 OpenVLA-OFT를 다음처럼 요약하면 부족하다.

\begin{warningbox}{Bad summary}
OpenVLA-OFT는 OpenVLA를 fine-tuning해서 성능을 크게 올린 논문이다. 최신 VLA fine-tuning 방법을 제안했고 LIBERO에서 좋은 결과를 보였다.
\end{warningbox}

같은 내용을 더 좋은 요약으로 바꾸면 다음과 같다.

\begin{sourcebox}{Good summary}
OpenVLA-OFT는 같은 OpenVLA base model을 유지한 상태에서 parallel decoding, action chunking, continuous action representation, L1 objective를 결합해 LIBERO 평균 success와 action generation throughput을 함께 개선한 fine-tuning recipe이며, 따라서 contribution은 새 backbone보다 action interface와 adaptation procedure에 있다 \citep{openvlaoft}.
\end{sourcebox}

좋은 요약은 짧아도 입력, 출력, 바뀐 mechanism, 검증 지표, 주장의 범위를 포함한다. 나쁜 요약은 ``좋다'', ``최신'', ``성능 향상''처럼 검증 조건을 지운다.

\section{리뷰와 연구 기획에서의 사용}

13개 질문은 읽기 도구이지만, 동시에 연구 기획 도구이다. 새 VLA 논문을 쓰기 전에 각 항목을 질문으로 뒤집으면 된다.

\begin{itemize}
  \item 배경: 독자가 왜 이 문제가 중요한지 바로 납득하는가?
  \item 문제: 과제가 너무 넓거나 모호하지 않은가?
  \item 기존 한계: prior work의 한계와 제안 방법이 직접 연결되는가?
  \item 목표: 목표가 측정 가능하고 검증 가능한가?
  \item 방법: 입력, 출력, 모듈, 학습/추론 절차가 재현 가능한가?
  \item 검증: 실험이 contribution을 직접 검증하는가?
  \item 비교: baseline이 충분히 강한가?
  \item 한계: reviewer가 찌를 약점을 선제적으로 다루는가?
\end{itemize}

이 방식은 논문 초안 평가에도 유용하다. 어떤 항목이 비어 있으면 그 부분이 reviewer의 공격 지점이 된다. 특히 VLA 분야에서는 benchmark와 baseline이 약하면 모델 구조가 좋아도 설득력이 급격히 떨어진다.

\section{요약}

VLA 논문을 읽을 때 13개 질문은 배경에서 공개 자원까지 주장을 닫는 도구이다. 이 프레임의 핵심은 방법을 멋지게 요약하는 것이 아니라, 입력과 출력, 학습 신호, 검증 조건, baseline, 실패 모드, 공개 자원을 같은 기준으로 확인하는 데 있다. 이후 장에서는 이 13개 질문을 각 주제별로 적용한다. Action representation 장에서는 질문 5와 6이 중심이 되고, benchmark 장에서는 질문 7--9가 중심이 되며, robustness 장에서는 질문 11과 12가 특히 중요해진다.

\begin{sourcebox}{Key papers for practicing the template}
\begin{itemize}
  \item Kim et al. (2024), OpenVLA: architecture, dataset mixture, fine-tuning, 공개 자원을 한 번에 점검하기 좋다.
  \item Kim, Finn, and Liang (2025), OpenVLA-OFT: 같은 base model에서 fine-tuning recipe가 claim을 어떻게 바꾸는지 보기 좋다.
  \item Pertsch et al. (2025), FAST: action tokenization claim과 closed-loop evidence를 분리해서 읽는 연습에 적합하다.
  \item Open X-Embodiment Collaboration (2023), Open X-Embodiment: dataset paper에서 claim, split, embodiment, source 공개를 확인하는 연습에 적합하다.
  \item Google DeepMind (2025/2026), Gemini Robotics page/report: closed industrial VLA에서 공개되지 않은 정보와 확인 가능한 정보를 구분하는 연습에 적합하다.
\end{itemize}
\end{sourcebox}

\begin{assignmentbox}{Reading assignment [Basic]}
같은 논문에 대해 13-question summary와 6-line quick note를 각각 작성하고, 두 형식 중 연구 주제 발굴에 더 도움이 된 항목이 무엇인지 설명하라.
\end{assignmentbox}
